\documentclass[11pt]{article}
\usepackage[a4paper,margin=18mm]{geometry}
\usepackage{amsmath,amssymb,amsthm,mathtools}

\usepackage[hidelinks]{hyperref}
\newcommand{\R}{\mathbb R}
\newcommand{\E}{\mathbb E}
\newcommand{\Id}{I}
\newcommand{\tr}{\operatorname{Tr}}
\newcommand{\Exp}{\operatorname{Exp}}
\newcommand{\Log}{\operatorname{Log}}
\newcommand{\proj}{\operatorname{proj}}

\newtheorem{proposition}{Proposition}
%\theoremstyle{remark}
\newtheorem*{remark}{Remark}

\title{Block spins and fluctuation integral}
\author{}
\date{}


\begin{document}
\maketitle
\abstract{The goal is to summarize the introduction of the block spin
transform from Section~3 and the fluctuation integral from Section~5.}

\setcounter{tocdepth}{1}
\tableofcontents

\newpage
\section{Gaussian block spin transforms}

This explains roughly Section~3.2.
%
The lattice $\Lambda$ has $n_{\rm fine}=|\Lambda|$ sites, $E$ oriented edges, and $B$ blocks.
Fields on the lattice are denoted $\phi$ and block spin fields $V$.

\subsection{Iterated Gaussian block spin transforms}

The initial kinetic term is given by the Laplacian on the lattice $\Lambda$:
\begin{equation}
  K_{\varnothing} = d^{\mathsf T} d.
\end{equation}
For a block $Y$, the linear block average defined by
\begin{equation}
 (Q\phi)_Y=\sum_{x\in Y}w_Y(x)\phi_x,
 \qquad w_Y(x)\ge0,\qquad \sum_{x\in Y}w_Y(x)=1.
\end{equation}
For a precision parameter $a>0$, the next-scale kinetic term is defined by
\begin{equation}
K_{h+1} = a\Id-a^2Q(K_h+aQ^*Q)^{-1}Q^*.
\end{equation}
For now, set $a=1$.

\begin{proposition}
For a given block spin field $V\in\R^B$, its quadratic form is the
minimum of the fine kinetic energy plus a penalty for disagreement
with the block average:
\begin{equation} \label{eq:Gauss-var}
 \frac12\langle V,K_{h+1}V\rangle
 =\min_{\phi\in\R^{n_{\rm fine}}} \mathcal E_V(\phi)
 =\min_{\phi\in\R^{n_{\rm fine}}}
 \left\{\frac12\langle\phi,K_h\phi\rangle
              +\frac a2\|V-Q\phi\|^2\right\}.
\end{equation}
Here $n_{\rm fine}$ is the number of fine sites and $B$ the number
of blocks.  The minimization is over unrestricted real fields,
and the mismatch is penalized rather than constrained to vanish.
\end{proposition}

\begin{proof}
Set $A=K_h+aQ^*Q$. For $a>0$, this matrix is positive definite whenever $\ker K_h\cap\ker Q=\{0\}$.
The positive gradient term in $K_h$ makes its kernel consist of constants, while
$Q$ preserves constants, so this condition holds.
Expanding the energy gives
\[
 \mathcal E_V(\phi)
 =\frac12\langle\phi,A\phi\rangle
  -a\langle\phi,Q^*V\rangle+\frac a2\|V\|^2.
\]
Its unique minimizer therefore satisfies
\[
 A\phi_V=aQ^*V,\qquad
 \phi_V=a(K_h+aQ^*Q)^{-1}Q^*V.
\]
Completing the square yields
\[
 \mathcal E_V(\phi)
 =\frac12\langle\phi-\phi_V,A(\phi-\phi_V)\rangle
  +\frac12\langle V,(a\Id-a^2QA^{-1}Q^*)V\rangle.
\]
The first term vanishes at $\phi=\phi_V$, giving the stated formula
for $K_{h+1}$.
\end{proof}


Note that $K_h$ is not finite range when $h \neq \varnothing$.
A sequence of block spin transformations is called \emph{history}
and the corresponding kinetic term $K_h$.


\subsection{Lifting the energy}

The renormalized kinetic term can be written as 
\begin{equation}
 K_h=c_0d^{\mathsf T}d+d^{\mathsf T}\ell_h^{\mathsf T}\ell_hd,
 \qquad
 \frac12\langle q,K_hq\rangle
 =\frac{c_0}{2}\|dq\|^2+\frac12\|\ell_hdq\|^2,
\end{equation}
where the matrix $\ell_h$ is indexed by edges:
\begin{equation}
  \ell_h : \R^E \to \R^E.
\end{equation}
There is redundancy between $c_0$ and $\ell_h$.
Initially, for any $c_0$,
\begin{equation}
 K_{\varnothing}=d^{\mathsf T}d,
 \qquad \ell_{\varnothing}=\sqrt{1-c_0}\,\Id.
\end{equation}
%
The constant $c_0>0$ is fixed (uniformly over histories) to reserve an explicit positive nearest-neighbour term for local energy estimates.
The operator $\ell_h$ has exponential tail.
Later $\ell_s$ is a spatially truncated version of $\ell_h$ used in the finite range integration.

\begin{proposition}
Define $D_0: \R^{n_{\rm fine}} \to \R^{E \times E \times B}$ by
\begin{equation}
  D_0
  =\begin{pmatrix}
       \sqrt{c_0}\,d\\ \ell_s d\\ -\sqrt a\,Q
  \end{pmatrix}
  =\begin{pmatrix}
       \mathsf A_s d\\ -\sqrt a\,Q
  \end{pmatrix},
 \qquad
 \mathsf A_s=\begin{pmatrix}\sqrt{c_0}\Id\\\ell_s\end{pmatrix}.
\end{equation}
%where
%\begin{equation}
% d,\ell_s d:\R^{n_{\rm fine}}\longrightarrow\R^E,
% \qquad Q:\R^{n_{\rm fine}}\longrightarrow\R^B.
%\end{equation}
Then the Hessian of the variational functional $\mathcal E_V(\phi)$ in \eqref{eq:Gauss-var} can be written as
\begin{equation}
 \|D_0\phi\|^2
 = \langle \phi, D_0^{\mathsf T} D_0 \phi \rangle
\end{equation}
and
\begin{equation}
  \mathcal E_V(\phi)- \mathcal E_V(\phi_V)
  = \frac12 \|D_0\phi +c_V - s_V\|^2,
\end{equation}
where $\phi_V$ is the minimizer of $\mathcal E_V$ and, for some $\nu_1,\nu_2,w$,
\begin{equation}
  c_V = \begin{pmatrix}
  0 \\ 0 \\ \sqrt{a}V
  \end{pmatrix},
  \qquad
  s_V
  = \begin{pmatrix}
    \nu_1 \\ \nu_2 \\ w
  \end{pmatrix}
  .
\end{equation}
\end{proposition}

\begin{proof}
The Hessian of the variational functional $\mathcal E_V(\phi)$ in \eqref{eq:Gauss-var} can be written as
\begin{equation}
 \|D_0\phi\|^2
 = \langle \phi, D_0^{\mathsf T} D_0 \phi \rangle
 =c_0\sum_e|(d\phi)_e|^2
  +\sum_e|(\ell_s d\phi)_e|^2
  +a\sum_Y|(Q\phi)_Y|^2
\end{equation}
and the full Gaussian energy is
\begin{equation}
  \mathcal E_V(\phi) = \frac12 \|D_0\phi + c_V\|^2 ,
  \qquad
  c_V = \begin{pmatrix}
    0 \\ 0 \\ \sqrt{a} V 
  \end{pmatrix}.
\end{equation}
Let $\phi_V = \operatorname{argmin} \mathcal E_V(\phi)$. Then 
\begin{equation}
  \mathcal E_V(\phi_V + \eta) = \frac12 \|D_0\eta + s_V\|^2  = \frac12 \|D_0\eta\|^2 + \frac12 \|s_V\|^2,
  \qquad
  s_V
  = D_0\phi_V + c_V
  = \begin{pmatrix}
    \nu_1 \\ \nu_2 \\ w
  \end{pmatrix}
\end{equation}
since $D^*s_V=0$. In other words,
\begin{equation}
  \mathcal E_V(\phi)- \mathcal E_V(\phi_V)
  = \frac12 \|D_0\phi - D_0\phi_V\|^2
  = \frac12 \|D_0\phi +c_V - s_V\|^2
\end{equation}
\end{proof}

\newpage
\section{Nonlinear background field coordinates}

This explains Section~5.2.

\subsection{Fluctuation integral}

Throughout this section, $S\subset \Lambda$ are the nonfrozen fine sites (small field region) with
 one spin $q_x \in S^2$ integrated at each site $x\in S$. 
Their number is $n=|S|$.
For each site $x\in S$, there is a background field of reference spins $T_x\in  S^2$
which are held fixed. We consider the integral
\begin{equation}
 \mathcal I=
 \int_{(S^2)^S} H(q)
    \exp\!\left\{-\frac{|R(q)|^2}{2g^2}\right\}
    \prod_{x\in S}d\omega(q_x),
 \label{eq:spin-integral}
\end{equation}
where $|R(q)|^2$ is a reference energy and $d\omega$ the normalized measure on $S^2$.
The function $H(q)$ contains nonlocal contributions not explicitly tracked in $|R(q)|^2$.


\subsection{Reference energy}

The precise form of the reference energy is unimportant for the algebraic manipulations
in the subsequent subsections. It is a nonlinear version of the Gaussian energy with soft block spin condition,
defined in the last section:
\begin{equation}
  \mathcal E_V(\phi) - \mathcal E_V(\phi_V) = \frac12 \|D_0\phi+c_V-s_V\|^2.
\end{equation}
More precisely, first restrict $D_0: \R^{n} \to \R^{E \times E \times B}$ 
to nonfrozen sites $S \subset \Lambda$ and `selected' components of the range $R\subset E\times E \times B$
(selecting small gradients or so):
\begin{equation}
D = D_0|_{S\times R}.
\end{equation}
%
Up to dropping rows and columns, in the last section it was shown that
\begin{equation}
D\phi +c_V -s_V = \begin{pmatrix}
  \sqrt{c_0} d\phi\\
  \ell_s d\phi\\
  \sqrt{a}V -\sqrt{a}Q\phi
\end{pmatrix}
-s_V,
  \qquad
  s_V = \begin{pmatrix}
    \nu_1 \\ \nu_2 \\ w
  \end{pmatrix}
.
\end{equation}
%
For $q: \R^n  \to S^2$, define $R(q)$ as a nonlinear version of this:
\begin{equation}
  R(q) = \begin{pmatrix}
  \sqrt{c_0} dq\\
  \ell_s dq\\
  \sqrt{a}(V-\widehat m(q))
  \end{pmatrix}
  - \begin{pmatrix}
    \nu_1 \\ \nu_2 \\ w
  \end{pmatrix}
\end{equation}
where 
\begin{equation}
  m_Y(q) = (Qq)_Y, \qquad \widehat m_Y(q) =
  \begin{cases}
      \frac{m_Y(q)}{|m_Y(q)|} & (|m_Y(q)|\geq \frac12)
      \\
      \dots
  \end{cases}
\end{equation}
is a nonlinear block spin average.

\subsection{Intrinsic tangent coordinates}

Consider a single site and reference spin $T \in S^2$.
For $u\in T^\perp$,
\begin{equation}
 \Exp_T u=\cos|u|\,T+\frac{\sin|u|}{|u|}u,
 \qquad \Exp_T0=T.
\end{equation}
The restriction $|u|<\pi$ parametrizes $S^2\setminus\{-T\}$ and the omitted point has area zero.
Note
\begin{equation}
 d\omega(\Exp_Tu)=J(u)\,d^2u,
 \qquad J(u)=\frac1{4\pi}\frac{\sin|u|}{|u|},
 \qquad J(0)=\frac1{4\pi}.
 \label{eq:area}
\end{equation}
In tangent coordinates, 
$q(u)=(\Exp_{T_x}u_x)_{x\in S}$ and
\begin{equation}
 \mathcal U=\prod_{x\in S}
       \{u_x\in T_x^\perp:|u_x|<\pi\}.
\end{equation}
The spherical area formula gives
\begin{equation}
 \mathcal I=
 \int_{\mathcal U} H(q(u))
 \exp\!\left\{-\frac{|R(q(u))|^2}{2g^2}\right\}
 \prod_{x\in S}\bigl[J(u_x)\,d^2u_x\bigr].
\end{equation}
Here $d^2u_x$ is Euclidean area measure on $T_x^\perp$.
This integral has $2n$ real coordinates.

\subsection{Auxiliary normal coordinate}

Introduce one real variable $v_x\in\R$ at each nonfrozen site.
Since $A=D^{\mathsf T}D$ is positive definite, the elementary
Gaussian integral gives
\begin{equation}
 1=\frac{(\det A)^{1/2}}{(2\pi g^2)^{n/2}}
       \int_{\R^n}\exp\!\left\{-\frac{v^{\mathsf T}Av}{2g^2}\right\}d^nv.
 \label{eq:scalar-one}
\end{equation}
Equivalently, the inserted probability law is a centered scalar
Gaussian with covariance $g^2A^{-1}$.  Its energy is
$v^{\mathsf T}Av=|Dv|^2$.  No cutoff is imposed on $v$, and
the inserted variables are initially independent of the spin variables.
Thus
\begin{equation}
 \begin{aligned}
 \mathcal I={}&\frac{(\det A)^{1/2}}{(2\pi g^2)^{n/2}}
 \int_{\R^n}\int_{\mathcal U} H(q(u))\\
 &\quad\times
 \exp\!\left\{-\frac{|R(q(u))|^2+|Dv|^2}{2g^2}\right\}
 \prod_{x\in S}\bigl[J(u_x)\,d^2u_x\, dv_x\bigr].
 \label{eq:uv-integral}
 \end{aligned}
\end{equation}
There are now $2n+n=3n$ real integration coordinates.
The purpose of the extra coordinate is
to complete each tangent plane to the ambient space via
\begin{equation}
\xi_x=u_x+T_xv_x
.
\end{equation}
Thus the scalar $v_x$ is an auxiliary normal coordinate, not a new physical spin component.

\subsection{Combining the coordinates into an ambient field}

Conversely, for $\xi=(\xi_x)_{x\in S}\in(\R^3)^S$, define the tangential and normal directions
\begin{equation}
 u_x(\xi)=\xi_x-(T_x\cdot\xi_x)T_x,
 \qquad
 v_x(\xi)=T_x\cdot\xi_x
\end{equation}
so that
\begin{equation}
 q_x(\xi)=\Exp_{T_x}u_x(\xi).
\end{equation}
The change of variables $(u_x,v_x) \mapsto \xi_x=u_x+T_xv_x$ has Jacobian one.
Thus
\begin{equation}
 \mathcal I=
 \frac{(\det A)^{1/2}}{(2\pi g^2)^{n/2}}
 \int_{\R^{3n}} H(q(\xi))\,\mathcal J(\xi)
 \exp\!\left\{-\frac{|R(q(\xi))|^2+|Dv(\xi)|^2}{2g^2}\right\}
 \,d^{3n}\xi ,
\end{equation}
where the chart factor is
\begin{equation}
 \mathcal J(\xi)=\prod_{x\in S} J(u_x(\xi)) 1_{|u_x(\xi)| < \pi}.
 %j(u)=
 %\begin{cases}
 %  \sin|u|/|u|,& |u|<\pi,\\
 %  0,& |u|\ge\pi,
 %\end{cases}
 %\qquad j(0)=1.
\end{equation}
%The factor $J(0)^n\mathcal J(\xi) = (4\pi)^{-n} \mathcal J(\xi)$ supplies the spherical
%Jacobian and the principal-chart restriction, while
%$(\det A)^{1/2}/(2\pi g^2)^{n/2}$ normalizes the auxiliary normal Gaussian.

\subsection{Gaussian coordinates}

Recall that $D: (\R^{3})^n \to (\R^3)^m$ acts on the nonfrozen coordinates.
Let $F=(F_i)_{i=1}^m\in(\R^3)^m$ be a fixed vector in the range so that
\begin{equation}
 (D\xi+F)_i=\sum_{x\in S}D_{ix}\xi_x+F_i.
\end{equation}
%
Define 
\begin{equation}
 \Delta(\xi)=|R(q(\xi))|^2+|Dv(\xi)|^2-|D\xi+F|^2.
 \label{eq:delta}
\end{equation}
Note that $R(q(\xi))$ is defined in terms of the physical sphere valued coordinates
whereas $D\xi+F$ is the affine ambient expression chosen for the Gaussian computation.
%
Clearly,
\begin{equation}
 e^{-(|R(q(\xi))|^2+|Dv|^2)/(2g^2)}
 =e^{-|D\xi+F|^2/(2g^2)}e^{-\Delta(\xi)/(2g^2)}.
\end{equation}
The first term on the right-hand side is a Gaussian density with
covariance $g^2A^{-1}$ where $A=D^{\mathsf T} D$.


\subsection{Gaussian expectation identity}

It will be useful to choose a Gaussian reference measure with more general covariance $g^2C$ than $g^2 A^{-1}$.
Choose any symmetric positive definite $n\times n$ matrix $C$
and any $\mu\in(\R^3)^S$.
Let $Z\in(\R^3)^S$ have the normalized Gaussian density with
\begin{equation}
\E[Z_x^\alpha Z_y^\beta]=C_{xy}\delta_{\alpha\beta}.
\end{equation}
%
The vector $\mu\in(\R^3)^S$ is an arbitrary fixed shift in
the Gaussian coordinates
\begin{equation}
 \xi=\mu+gZ,\qquad Z\sim N(0,C\otimes I_3).
\end{equation}
The covariance $C$ has not been determined yet, but it will have finite range.

\begin{remark}
The unique minimizer of the Gaussian action
\begin{equation}
\mathcal E_F(\xi)=
\frac12|D\xi+F|^2
\end{equation}
is 
\begin{equation}
\xi = \mu_*=-A^{-1}D^{\mathsf T}F.
\end{equation}
\end{remark}

\begin{proof}
\begin{equation}
 \nabla_\xi\frac12|D\xi+F|^2=A\xi+D^{\mathsf T}F,
\end{equation}
so positive definiteness of $A$ makes $\mu_*$ the unique minimizer. 
\end{proof}

This above choice $\mu=\mu_*$ gives $h=0$ in the identity below, eliminating the term linear in $Z$ in the exponent in the following proposition.
The general identity below permits other choices. The preprint uses a finite range approximation
\begin{equation}
 \mu_s=-\Pi_0^{\mathsf T}D^{\mathsf T}F,
\end{equation}
related to the choice $C=C_s$ where
\begin{equation}
  C_s = \frac12 (\Pi_0+\Pi_0^{\mathsf T}),
\end{equation}
where the finite range condition is enforced by
\begin{equation}
  (\Pi_0)_{xy} = 1_{y\in\Omega_x} (A_{\Omega_x} +a_0)^{-1}_{xy}.
\end{equation}
For this choice, $h=A\mu_s+D^{\mathsf T}F$ need not vanish;
its contribution is retained explicitly.

\begin{proposition}
Given arbirary $F$, for any $C$ and $\mu$,
\begin{align}
 \mathcal I={}&
 \kappa_g^n(\det A)^{1/2}(\det C)^{3/2}
 \exp\!\left\{-\frac{|r|^2}{2g^2}\right\}
 \nonumber\\
 &\E_{\gamma_C}\left[
 H(q(\mu+gZ))\,\mathcal J(\mu+gZ)
 \exp\!\left\{
 -\frac{\Delta(\mu+gZ)}{2g^2}
 +\frac12 Z^{\mathsf T}(C^{-1}-A)Z
 -g^{-1}(h,Z)
 \right\}\right],
 \label{eq:main}
\end{align}
where %$h^{\mathsf T}Z=\sum_xh_x\cdot Z_x$ and
\begin{equation}
 r=D\mu+F,\qquad h=A\mu+D^{\mathsf T}F,
 \qquad \kappa_g=(2\pi g^2)J(0)=\frac{g^2}{2}.
\end{equation}
\end{proposition}

\begin{proof}
By the above,
\begin{align}
 \mathcal I={}&
 \frac{(\det A)^{1/2}}{(2\pi g^2)^{n/2}}
 \int_{\R^{3n}}H(q(\xi))\,\mathcal J(\xi)
 \exp\!\left\{-\frac{|D\xi+F|^2+\Delta(\xi)}{2g^2}\right\}d^{3n}\xi.
 \label{eq:lebesgue}
\end{align}
Now set $\xi=\mu+gZ$, whose Jacobian is $g^{3n}$.
Ordinary multiplication of squares gives
\begin{equation}
 |D(\mu+gZ)+F|^2
 =|r|^2+2g\,h^{\mathsf T}Z+g^2Z^{\mathsf T}AZ.
\end{equation}
The constant is
\begin{equation}
 \frac{(\det A)^{1/2}}{(2\pi g^2)^{n/2}}
 g^{3n}(2\pi)^{3n/2}(\det C)^{3/2}
 =\kappa_g^n(\det A)^{1/2}(\det C)^{3/2}.
\end{equation}
\end{proof}


%\newpage
%\section{Parametrization of the renormalized action}

%The renormalized action or rather its exponential (called `retained density') is parametrized in (3.29).

\newpage
\section{Finite range fluctuation integrals}

One of the next steps to understand seems to be Section 5.4 of the first O4 paper.

\subsection{Truncation of the Gaussian covariance}

The finite range is enforced by the condition $y \in \Omega_x$. The set $\Omega_x$ is referred to as \emph{inverse window}.
It is defined by
\begin{equation}
 .... 
\end{equation}
Therefore
\begin{equation}
C^{-1}-A = ...
\end{equation}


\end{document}
